% !TEX root = userguide.tex

\def\headdate{6th February 2009}

\section{Viscoelastic fluid flow using stochastic models (BCF)}
\label{sec:bcfdg}

\subsection{Weak form of the equations}

We assume the user is already familiar with Secs.~\ref{sec:weakviscoelastic}
and \ref{sec:DGweak}.

The equations for the flow of a single-mode viscoelastic fluid
in a region $\Omega$ are given by Eqs.~\eqref{eq:ve1}--\eqref{eq:ve3}. For the
BCF computations the viscoelastic fluid model is replaced by the BCF equation
(derived from Eq.~\eqref{eq:Langevin}) \cite{Hulsen97}:
\begin{equation}
  \mathrm{d}\vek{Q}=
     (-\vek u\cdot\nabla\vek Q+\ten L\cdot\vek Q - \frac{1}{2\lambda_H}\vek F_c)
    \mathrm{d}t+\frac{1}{\sqrt{\lambda_H}}\mathrm{d}\vek W,
\end{equation}
where $\vek Q(\vek x,t)$ is a configuration field and
$\mathrm{d}\vek W(t)$ is the Wiener increment. The stress tensor
 $\ten\tau$ is now
determined by the ensemble $\vek Q_\text{ens}$ using the usual stress expression
Eq.~\eqref{eq:sttress} extended to the full domain $\Omega$.

For the discontinuous Galerkin \fem\ we first start with a finite element mesh
before defining the weak form:
\[
\Omega=\cup_e\Omega_e, \quad\Omega_{e}\cap\Omega_{f}=\emptyset\quad
   \text{for }e\neq f
\]
The weak form can be now be stated, element wise for the BCF equation,
as:
Find $\vek u$, $p$ and $\ten c$ such that
\begin{alignat}{2}
 \big(\vek v,\rho(\pderiv{\vek u}t+\vek u\cdot\nabla\vek u)\big) +
  \big((\nabla\vek v)^T,2\eta_s\ten D\big)
    &-(\nabla\cdot\vek v,p)+\big((\nabla\vek v)^T,\ten\tau(\vek
Q_\text{ens}\big)\notag\\
    &=(\vek v,\vek t_{\text{N}})_{\Gamma_{\text{N}}}+(\vek v,\rho\vek b)
&&\text{for all }\vek v\\
  \,(q,\nabla\cdot\vek u)&=0,&&\text{for all }q\\
 \int_{\Omega_e}\ten R:\big(
 \mathrm{d}\vek Q +\vek u\cdot\nabla\vek Q-\ten L\cdot\vek Q + \frac{1}{2\lambda_H}\vek F_c)
    \mathrm{d}t&-\frac{1}{\sqrt{\lambda_H}}\mathrm{d}\vek W,
\big)\D x\notag\\
    &+\int_{\Gamma^e_{\text{in}}}(\vek u\cdot\vek n)\vek R:[\vek Q]\D s = 0,
  &&\text{for all }\vek R,\Omega_e
   \label{eq:weakQDG}
\end{alignat}
using appropriate spaces for $\vek u$, $p$, $\vek Q$, $\vek v$, $q$ and 
$\vek R$.
The function space for $\vek Q$ and $\vek R$ is discontinuous on the
element boundaries.
In equation~\eqref{eq:weakQDG} $\Gamma^e_{\text{in}}$ is the
inflow boundary of element $e$ and
$[\vek Q]=\vek Q^+-\vek Q$ is the ``jump'' across that element boundary.
Also, $\vek n$ is the outwardly directed unit normal vector on the boundary of
the element and $\vek Q^+$ is
the `outside' value of $\vek Q$.

The weak form can be used to obtain an approximate solution
by interpolating $\vek u$, $p$ and $\vek Q$ 
(and similarly $\vek v$, $q$ and $\vek R$) by
polynomials on each element:
\begin{align}
    \vek u_h&=\tsum_k \phi_k(\vek x)\vek u_k=\col\phi^T(\vek x)\col{\vek u}\\
    p_h&=\tsum_k \psi_k(\vek x) p_k=\col\psi^T(\vek x)\col{ p} \\
  \ten Q_h&=\tsum_k \theta_k(\vek x)\vek Q_k=\col\theta^T(\vek x)\col{\vek Q}
  \label{eq:thetaQDG}
\end{align}
where $\col\phi(\vek x)$, $\col\psi(\vek x)$, $\col\theta(\vek x)$
are global shape functions. Note, that most terms in the momentum balance and
the continuity equation can be build using standard (Navier-) Stokes elements.

As explained in Sec.~\ref{sec:timeDG} the `mass' matrix $\mat M$ of the
discretized weak form of Eq.~\eqref{eq:weakQDG} is block diagonal:
\begin{equation}
   \mat M=\begin{pmatrix}
              \mat M_1 & \mat 0 & \hdotsfor{2} & \mat 0 \\
              \mat 0 & \mat M_2 & \mat 0 & \ldots & \mat 0 \\
                              \vdots &  && \ddots  & \vdots \\
              \mat 0 & \hdotsfor{2} & \mat 0 & \mat M_{N_e}
      
          \end{pmatrix}
\end{equation}
with $\mat M_e$ the mass matrix of element $e$. We can take this one step
further to a fully diagonal matrix by
\begin{itemize}
   \item using a Gauss rule for integration of the volume integrals having the
     same number of points as the number of `nodes' or unknowns per element for
     a single component of $\vek Q$,
   \item locating the nodes of the element in the Gauss points.
\end{itemize}
This procedure is called collocation. For example with $Q_2-P_1^d$
 velocity-pressure
elements and $Q_1$ interpolation for the $\vek Q$,
we have four nodes for $\vek Q$
in each element. Using a 2$\times$2 Gauss-Legendre integration rule and
locating the nodes in the Gauss points will lead to collocation. We still have
to use the 3$\times$3 Gauss integration
for the momentum balance for accuracy reasons. For the collocation procedure
we split the volume integral of the weak form of the Langevin equation 
Eq.~\eqref{eq:weakQDG} into two parts:
\begin{multline}
 \int_{\Omega_e}\ten R:\Big(
 (\mathrm{d}\vek Q -\ten L\cdot\vek Q + \frac{1}{2\lambda_H}\vek F_c)
    \mathrm{d}t-\frac{1}{\sqrt{\lambda_H}}\mathrm{d}\vek W
\Big)\D x\\+
 \int_{\Omega_e}\ten R:\big(
 \vek u\cdot\nabla\vek Q \big)\D x
    +\int_{\Gamma^e_{\text{in}}}(\vek u\cdot\vek n)\vek R:[\vek Q]\D s = 0,
  \label{eq:weakQDG2}
\end{multline}
If we apply the collocation procedure the coupling between the nodes/Gauss
points for $\vek Q$ fully disappears for the first part. Multiplying the
equation with the inverse of the diagonal mass matrix shows that the first part
on integration point level is identical to what we have in a traditional
Brownian dynamics simulation. There is an extra term due to the second
part (convection of fields $\vek Q$). Note, that we do not fully apply the
collocation procedure for the volume integral of the convection term
(the second term): we use the $3\times3$ Gauss integration rule used for the
momentum balance (for $Q_2-P_1^d$), but the
shape functions of $\vek Q$ are still based on the nodes in the $2\times2$
Gauss points (for $Q_1$), of course.

We will be using the DEVSS formulation for the momentum balance
(see Sec.~\ref{sec:weakviscoelastic}).

\subsection{Time discretization (without inertia)}

Assuming $\vek u^{n}$ and $\vek Q^n_\text{ens}$ is known we can compute
$\vek Q^{n+1}_\text{ens}$ by
time advancing Eq.~\eqref{eq:weakQDG2} using methods from Brownian dynamics
treating the convection term as an explicit force.
Then we compute $\vek u^{n+1}$, 
$\vek p^{n+1}$ and $\ten E^{n+1}$ from the momentum balance,
continuity equation and projection equation at time $t^{n+1}$:
\begin{align}
      -\nabla\cdot\big(2\eta_s\ten D(\vek u^{n+1})\big) +\nabla p^{n+1} 
  -2\alpha\nabla\cdot\big( \ten D(\vek u^{n+1}) -\ten E^{n+1} \big)
&=\nabla\cdot\ten \tau(\vek Q^{n+1}_\text{ens})
        + \rho\vek b \label{eq:mombalQDG}\\
     \nabla\cdot\vek u^{n+1}&=0\\
   -\ten D(\vek u^{n+1})+\ten E^{n+1} &= \ten 0
\end{align}
We call this scheme the \emph{explicit stress} formulation: the stress tensor
$\ten \tau(\vek Q^{n+1}_\text{ens})$ in the momentum balance is
computed in an
explicit way from $\vek Q^{n+1}_\text{ens}$, which itself has been
computed first from the Langevin equation.

The time integration procedure described above is only suitable for fluids
where $\eta_s\neq0$. For fluids with $\eta_s=0$, such as melts described by a
large number of relaxation times, other procedures are required.
For small values of
$\eta_s$ the time step limitation can be very severe and the explicit stress
formulation cannot be used either. 

\subsection{Planar Couette flow of a Hookean dumbbell}
\label{sec:bcfcouette}

We consider a Hookean dumbbell fluid for a planar Couette problem on a
unit square.
The problem is similar to the problem for an Oldroyd-B fluid in
Sec.~\ref{sec:couette}, but now using the BCF/DG method.
The problem is part of the example files in `bcf':
\begin{quote}
   \texttt{getexample bcf}
\end{quote}
There is one example file: \texttt{couette1.f90}.
The input file is \texttt{input\_couette1.txt}.
The reader is advised to study the example in detail.
For the equation \eqref{eq:weakQDG2} it is not necessary to build a matrix, nor
is it necessary to build the vector by FEM assembly, because the unknowns on
element level are are decoupled from other elements. Therefore we use a
simplified `build-system' called \texttt{loop\_over\_elements} in which the
system vector is directly written to a vector in the structure
\texttt{oldvectors\_t}.

\subsection{Flow around a cylinder between two plates for a Hookean dumbbell
model.}

The problem is part of the example files in `bcf':
\begin{quote}
   \texttt{getexample bcf}
\end{quote}
There is a mesh-generation file: \texttt{mesh\_confined\_cylinder.f90}.
There are two main files: \texttt{cylinder1.f90} and \texttt{cylinder2.f90}.
The input files are \texttt{input\_cylinder1.txt} and
\texttt{input\_cylinder2.txt}, respectively. The main program
\texttt{cylinder2.f90} differs only from \texttt{cylinder1.f90} by the addition
of the output of $\det\ten c$ in various ways.
The post-processing files are \texttt{post\_cylinder.f90} and 
\texttt{streamfunction.f90}.

NOTE: there used to be an example called \texttt{bcf\_sepmesh}, where the same
cylinder problem is given, however using a \sep-generated mesh. 
This example has been moved to the legacy repository.

\subsection{Flow around a cylinder between two plates for a FENE dumbbell
model.}

The problem is part of the example files in `bcf':
\begin{quote}
   \texttt{getexample bcf}
\end{quote}
There is a mesh-generation file: \texttt{mesh\_confined\_cylinder.f90}.
The main file is: \texttt{cylinder3.f90}.
The input file is \texttt{input\_cylinder3.txt}.
The post-processing files are \texttt{post\_cylinder.f90} and 
\texttt{streamfunction.f90}.

